% Source-derived chapter 12: Matching of sheaves
% Original source: higher-siegel-weil-unitary-nonsingular-terms
\chapter{Matching of sheaves}
\label{higher-siegel-weil-unitary-nonsingular-terms-chapter-12}
\source{higher-siegel-weil-unitary-nonsingular-terms}

\begin{remark}[Source section]
\label{higher-siegel-weil-unitary-nonsingular-terms-chapter-12-source}
\source{higher-siegel-weil-unitary-nonsingular-terms}
\source{higher-siegel-weil-unitary-nonsingular-terms}
This chapter reproduces the corresponding section of the retrieved
LaTeX source, including its definitions, statements, examples, and
proofs. It has not yet been translated into Lean.
\notready
\end{remark}

% The source section title was: Matching of sheaves
\subsection{Recap} 
Let 
\begin{equation}\label{eq: normalized Eisenstein coefficients}
\source{higher-siegel-weil-unitary-nonsingular-terms}
\wt{E}_{a}(m(\cE),s,\Phi)   = E_a(m(\cE), s, \Phi) \cdot \chi(\det(\cE))^{-1} q^{\deg(\cE)( s-\frac{n}{2})+\frac{1}{2}n^2\deg(\omega_X) } \cdot  \sL_n(s) = \Den(q^{-2s}, (\cE,a))
\end{equation}
where the notation is as in Theorem \ref{th:Eis Den}, being a renormalization of the $a^{\mrm{th}}$ Fourier coefficient of $E_a(m(\cE), s, \Phi)$.

We emphasize that, in keeping with \S \ref{ssec: notation}, $X$ is proper and $\nu \co X' \rightarrow X$ is a finite \'{e}tale double cover (possibly trivial). 

\begin{thm}
\source{higher-siegel-weil-unitary-nonsingular-terms}
\notready\label{thm: main} Keep the notations above. Let $(\cE,a)\in \cA_{d}(k)$. Then we have 
\source{higher-siegel-weil-unitary-nonsingular-terms}
\begin{equation}\label{eq: SWF}
\source{higher-siegel-weil-unitary-nonsingular-terms}
\deg [\cZ_{\cE}^r(a)] = \frac{1}{(\log q)^r}  \left(\frac{d}{ds}\right)^r\Big|_{s=0} \left( q^{ds} \wt{E}_{a}(m(\cE),s,\Phi) \right). 
\end{equation}
\end{thm}

In the previous parts, we have found sheaves on $\cA_d$ which correspond to the two sides of \eqref{eq: SWF}, in the sense of the function-sheaf dictionary. Let us summarize the situation. 

 On the analytic side, we proved a formula expressing the non-singular Fourier coefficient of the Siegel--Eisenstein series in terms of the Frobenius trace of a graded virtual perverse sheaf $\cK_{d}^{\Eis}(T)$ on $\Herm_{2d}(X'/X)$.

\begin{thm}[Combination of Theorems \ref{th:Eis Den} and \ref{th:Den Wd}]
\source{higher-siegel-weil-unitary-nonsingular-terms}
\notready Let $(\cE,a)\in \cA_{d}(k)$. Then we have 
\begin{equation}\label{eq: comparison 1}
\source{higher-siegel-weil-unitary-nonsingular-terms}
\wt{E}_{a}(m(\cE),s,\Phi) =  \Tr(\Frob,\cK_{d}^{\Eis}(q^{-2s})_{\cQ}).
\end{equation}
\end{thm}


On the geometric side,  in Corollary \ref{cor: KR degree as trace of Frob}, we found a formula expressing the degree of the special $0$-cycle in terms of $r^{\mrm{th}}$ derivative of the Frobenius trace of another graded perverse sheaf $\cK^{\Int}_{d}(T)$ on $\Herm_{2d}(X'/X)$, repeated below: 
\begin{equation}\label{eq: comparison 2}
\source{higher-siegel-weil-unitary-nonsingular-terms}
\deg [\Cal{Z}^{r}_{\Cal{E}}(a)] = \frac{1}{(\log q)^r} \left(\frac{d}{ds}\right)^r\Big|_{s=0} \left(q^{ds}\Tr(\Frob, \cK_d^{\mrm{Int}}(q^{-2s})_{\cQ})\right).
\end{equation}


\subsection{Proof of the main theorem} Comparing \eqref{eq: comparison 1} and \eqref{eq: comparison 2}, we see that in order to prove Theorem \ref{thm: main}, it remains to match the graded sheaves $\cK_d^{\mrm{Int}}(T)$ and $\cK_{d}^{\Eis}(T)$ on $\Herm_{2d}(X'/X)$.

\begin{prop}
\source{higher-siegel-weil-unitary-nonsingular-terms}
\notready
We have $ \Cal{K}^{\Int}_{d}(T) \cong \cK_{d}^{\Eis}(T)$ as graded perverse sheaves  on $\Herm_{2d}(X'/X)$.
\end{prop}
\begin{proof} Both sides can be written as $\Spr^{\Herm}_{2d}[\r]$ for some graded virtual representation $\r$ of $W_{d}$. By definition (Definition \ref{def:Kd}), the sheaf $\cK^{\Int}_{d}(T)$ corresponds to 
\[
\r_{\cK_{d}^{\Int}}(T)=\sum_{i=0}^{d} \Ind_{W_i \times W_{d-i}}^{W_d} (\chi_i \boxtimes 1) T^{i}.
\]
We calculate the (a priori virtual) representation of $W_d$ which corresponds under Springer theory to the $\cK_{d}^{\Eis}(T)$ from Definition \ref{def: K_Eis}. The operation $R\oll{f}_{i,!}R\orr{f}^{*}_{i}$ corresponds under Springer theory to $\Ind_{W_{d-i} \times S_i}^{W_d}$ and $\frP_{d-i}(T)$ corresponds to $\sum_{j=0}^{d-i} (-1)^j \Ind_{W_j \times W_{d-i-j}}^{W_{d-i}} (\sgn_j \boxtimes \mbf{1}) T^j$ (cf. Definition \ref{d:Spr isotypic}  for $\HSpr_{d}$). Hence $
\cK_{d}^{\Eis}(T)$ corresponds to 
\[
\sum_{i=0}^d \Ind_{W_{d-i} \times S_i}^{W_d} \left( \sum_{j=0}^{d-i}(-1)^j  \Ind_{W_j \times W_{d-i-j}}^{W_{d-i}}(\sgn_j \boxtimes \one_{d-i-j}) \boxtimes \one_i T^{i+j} \right) .
\]
After re-indexing, we have
\[
\r_{\cK_{d}^{\Eis}}(T)=\sum_{i=0}^{d} \sum_{j=0}^{i}(-1)^j\Ind_{S_{i-j} \times W_j \times W_{d-i}}^{W_d} (\one \bt \ol{\sgn}_{j} \bt \one ) T^{i}.
\]
The desired statement then follows from the Lemma below (whose notation has been re-indexed) by comparing each coefficient.
\end{proof}


\begin{lemma}
\source{higher-siegel-weil-unitary-nonsingular-terms}
\notready We have the identity of virtual representations of $W_d$: 
\[
\chi_d = \sum_{j=0}^{d} (-1)^j \Ind_{S_{d-j} \times W_j}^{W_d}(1 \boxtimes \ol{\sgn}_{j}).
\]
\end{lemma}
\begin{proof}
We will prove this by comparing traces of an arbitrary element $g \in W_d$.  For $g \in W_d$, 
\begin{equation}\label{eq: trace of induced}
\source{higher-siegel-weil-unitary-nonsingular-terms}
\Tr(g, \Ind_{S_{d-j} \times W_j}^{W_d} (1 \times \ol{\sgn}_j))=\sum_{\substack{w\in W_{d}/(S_{d-j} \times W_j) \\ w^{-1}gw\in S_{d-j} \times W_j} } \ol{\sgn}_j(g'').
\end{equation}
Here, when $w^{-1}gw\in S_{d-j} \times W_j$, we write $w^{-1}gw=(g',g'')$ for $g'\in S_{d-j}$ and $g''\in W_{j}$. 

Identify $W_d$ with the group of permutations of $\{\pm 1, \ldots, \pm d\}$ that commute with the involution $\s$ exchanging $j \leftrightarrow -j$ for all $1 \leq j \leq d$. The subgroup $S_{d-j}\times W_{j}$ is the stabilizer of $\{1,2,\cdots, d-j\}$. Therefore the coset space $W_{d}/(S_{d-j} \times W_j)$ is in natural bijection with subsets $J\subset \{\pm 1, \ldots, \pm d\}$ such that $|J|=d-j$ and $J\cap (-J)=\vn$. Let $\frJ_{g}$ be the set of $J\subset \{\pm 1, \ldots, \pm d\}$ such that $|J|=d-j$, $J\cap (-J)=\vn$ and $gJ=J$.  Let $g''_{J}$ be the permutation of $g$ on $\{\pm 1, \ldots, \pm d\}\bs (J\cup (-J))$. Combining this with \eqref{eq: trace of induced}, we obtain 
\begin{equation}
\sum_{j=0}^{d}(-1)^{j}\Tr(g, \Ind_{S_{d-j} \times W_j}^{W_d} (1 \times \ol{\sgn}_j))=\sum_{J\in \frJ_{g}}(-1)^{d-|J|}\ov\sgn(g''_{J}).
\end{equation}




%We have an embedding $S_d \hookrightarrow W_d$ whose image consists of $g\in W_{d}$ that stabilizes $\{1,2,\cdots, d\}$. 

%is a sum over cosets $w(S_{d-j} \times W_j)\in W_{d}/(S_{d-j} \times W_j)$ suall ways to conjugate $g$ to $\underbrace{g'}_{\in S_{d-j}}\underbrace{g''}_{\in W_j}$ of $\ol{\sgn}_j(g'')$. 

%We first explicate the conjugacy classes in $W_d$. Identify $W_d$ with permutations of $\{\pm 1, \ldots, \pm d\}$ that commutes with the involution $w_0$ exchanging $j \leftrightarrow -j$ for all $1 \leq j \leq d$. We have an embedding $S_d \hookrightarrow W_d$ whose image consists of $g\in W_{d}$ that stabilizes $\{1,2,\cdots, d\}$.

For any $g \in W_d$,  the cycle decomposition of $g$ can be grouped into a decomposition $g = g_1 \ldots g_r$ (unique up to reordering) where $g_i$ is one of the two forms:
\begin{itemize}
\item (positive bicycle) $g_i$ is a product of two disjoint cycles $c_{i}\s(c_{i})$ (in particular, no two elements appearing in $c_{i}$ are negatives of each other). 
\item (negative cycle) $g_i$ is a single cycle invariant under the involution $\s$. 
\end{itemize}
Let $C^{+}_{g}$ be the set of cycles of $g$ that are part of a positive bicycle (i.e., $C_g^+$ contains both $c_{i}$ and $\s(c_{i})$ for each positive bicycle $g_{i}$). For any $x\in W_{d}$ we denote by $\un x\subset \{\pm1,\cdots, \pm d\}$ the set of elements that are not fixed by $x$. For a cycle $c$ we let $|c|$ be its length. From this description we see that $J\in \frJ_{g}$ if and only if $J$ is a union of $\un c$ for a subset of cycles  $c\in C^{+}_{g}$. In other words, consider the set $\frI_{g}$ of subsets $I\subset C^{+}_{g}$ such that $\un I$ is disjoint from $\s(\un I)$. Then we have a bijection $\frI_{g}\isom \frJ_{g}$ sending $I\in \frI_{g}$ to $J:=\cup_{c\in I}\un c$.

For $I\in \frI_{g}$, let $g'_{I}$ be the product of  $g_{i}$ such that $g_{i}$ contains a cycle in common with $I$; let $g''_{I}$ be the product of the remaining $g_{i}$'s. The above discussion allows us to rewrite
\begin{equation}
\sum_{j=0}^{d}(-1)^{j}\Tr(g, \Ind_{S_{d-j} \times W_j}^{W_d} (1 \times \ol{\sgn}_j))=\sum_{I\in \frI_{g}}(-1)^{\sum_{c\in C^{+}_{g}}|c|}\ov\sgn(g''_{I}).
\end{equation}

This sum factorizes as a product over the $g_i$ with individual factors as follows: %, over $I \subset \{ g_1, \ldots, g_r\}$ indexing cycles of Type 1 with $|I| = n-j$, of the  $\prod_{i \in I} g_i$ into $S_{n-j}$ times $\sgn$ applied to the projection of the remaining cycles to $S_j$. 
\begin{itemize}
\item For a positive bicycle $g_i = c_i \s(c_{i})$, the local factor is the sum of three contributions, corresponding to whether $c_{i}\in I$, $\s(c_{i})\in I$ or neither  $c_{i}$ nor $\s(c_{i})$ is in $I$. The first two cases each contribute $1$. The last case leads to a contribution of $(-1)^{|c_i|} \sgn(c_i)=-1$. The total contribution of the factor corresponding to a positive bicycle $g_{i}$ is therefore $1+1+(-1)=1$.

\item For each negative cycle $g_i$, since it always appears in $g''_{I}$, its contribution is  $(-1)^{|g_{i}|/2}\ol{\sgn}(g_i)$. Let  $\ov g_{i}$ be the image of $g_{i}$ in $S_{d}$, which is a cycle of length $|g_{i}|/2$. Then $(-1)^{|g_{i}|/2}\ol{\sgn}(g_i)=(-1)^{|\ov g_{i}|}\sgn(\ov g_{i})=-1$. Therefore the contribution of the factor corresponding to a negative cycle $g_{i}$ is $-1$.
%We can express $g_i$ as a product of transpositions of the form $(j, \,  -j)$ (appearing at most once for each $j$) followed by a positive bicycle  $(x_1, \, \ldots \, , x_k)(-x_1,  \, \ldots \, , -x_k)$ where each $j$ appears among the $x_k$. Then $(-1)^k \ol{\sgn}_k(g_i) = -1$. 
\end{itemize}
Summarizing, we have found 
\begin{equation}\label{Trg factor}
\source{higher-siegel-weil-unitary-nonsingular-terms}
\sum_{j=0}^{d} (-1)^{j}\Tr(g, \Ind_{S_{d-j} \times W_j}^{W_d} (1 \times \ol{\sgn}_j)) = \left( \prod_{g_i \text{ positive }} 1 \right) \cdot  \left( \prod_{g_i \text{ negative} }(-1) \right).
\end{equation}

%\underbrace{(1)}_{=\chi_n(g_i)}
On the other hand,  we have
\begin{equation}
\chi_{d}(g_{i})=\begin{cases} 1, & g_{i} \mbox{ is positive;} \\ -1, &g_{i} \mbox{ is negative.} \end{cases}
\end{equation}
Indeed, if $g_{i}=c_{i}\s(c_{i})$ is positive, then we have $\chi_{d}(c_{i})=\chi_{d}(\s(c_{i}))=1$ because both $c_{i}$ and $\s(c_{i})$ can be conjugated into $S_{d}$. If $g_{i}$ is negative, then up to conjugacy we may assume $g_{i}$ is the cyclic permutation $(1,2,\cdots, m, -1,\cdots, -m)$ for some $1\le m\le d$. Then $g_{i}=(1,-1)(1,2,\cdots, m) (-1,-2,\cdots, -m)$, from which we see $\chi_{d}(g_{i})=-1$.

We conclude that the right side of \eqref{Trg factor} is $\prod_{}\chi_{d}(g_{i})=\chi_{d}(g)$. This completes the proof.
%i.e. when we write $g_i$ in the form above there are an odd number of transpositions of the form $(j \, -j)$. Indeed, it is straightforward to verify that if there are an even number of transpositions, then $g_i$ would break into a product of two disjoint cycles. 


%Two elements $g,g' \in W_d$ are conjugate if and only if their factors $\{g_i\}$ and $\{g_i'\}$ are equivalent under some permutation of $\{1, \ldots, d \}$. Hence $g$ to conjugate to $g'g'' \in S_{d-j} \times W_j$ if and only if there is a subset of its $\{g_i\}$ involving $d-j$ elements in $\{1, \ldots, d\}$. Say this subset is indexed by the set $I$. Then the number of ways to conjugate $g$ into $S_{d-j} \times W_j$ is $2^{|I|} \# (S_{d-j} \times W_j)$, with the factor of $2^{|I|}$ coming from the fact that $c_i$ is conjugate to $c^{-}_i$ in $W_d$.  \mnote{does this require more explanation?} 
%
%Therefore, $\sum_j \Tr(g, \Ind_{S_{d-j} \times W_j}^{W_d} (1 \times \ol{\sgn}_j))$ is the sum over all subsets $I \subset \{ g_1, \ldots, g_r\}$ indexing cycles of Type 1 with $|I| = d-j$, of $2^{d-j}$ times $\sgn$ applied to the projection of the remaining cycles to $S_j$. 

\end{proof}


\subsection{The split case $X'=X\coprod X$}
We make our result more explicit in the split case $X'=X\coprod X$.
  
  On the analytic side in \S\ref{ss:Eis}, the group $H_n=\GL_{2n,F}$ and $P_n$ is the standard parabolic corresponding to the partition $(n,n)$, with Levi $M_n\simeq \GL_{n,F}\times \GL_{n,F}$. We then have the degenerate principal series
  $$
I_n(s)=\Ind_{P_n(\BA)}^{H_n(\BA)}( |\cdot|_{F}^{s+n/2}\times |\cdot|_{F}^{-s-n/2}),\quad s\in \mathbb{C}.
$$
Let $\cE=(\cE_1,\cE_2)\in \Bun_{M_n}(k)\simeq\Bun_{\GL_n}(k)\times \Bun_{\GL_n}(k)$, and let $a:\cE_1\to\cE_2^\vee$ be an injective map of $\cO_X$-modules. Then by \S\ref{ss:FC} the Siegel--Eisenstein series has a well-defined $a^{\mrm{th}}$ Fourier coefficient $E_{a}(m(\cE),s,\Phi)$ at $(\cE_1,\cE_2)$. By Theorem \ref{th:Eis Den}  and \ref{th:Eis Den torsion} we have
  \begin{eqnarray*}
E_{a}(m(\cE),s,\Phi)=q^{-(\deg(\cE_1)+\deg(\cE_2))( s-n/2)-\frac{1}{2}n^2\deg\omega_X }\sL_n(s) ^{-1} \Den(q^{-2s},  \cE_2^\vee/\cE_1),
\end{eqnarray*}
where $\sL_n(s)=\prod_{i=1}^n\zeta_F(i+2s)$ and, for a torsion $\cO_X$-module $\cQ$, the density polynomial is given by
$$
 \Den(T,  \cQ)=\sum_{0 \subset \cI_1 \subset \cI_2\subset \cQ} T^{\dim_{k}\cI_1+\dim_k \cQ/\cI_2}\prod_{v\in |X|}\fm_{v}(t_v(\cI_2/\cI_1);T^{\deg(v)}).
$$
Here see \eqref{def:maT} for $\fm_{v}(t_v;T)$.
The normalized Fourier coefficient \eqref{eq: normalized Eisenstein coefficients} is
\[
\wt{E}_{a}(m(\cE),s,\Phi)  =  \Den(q^{-2s},  \cE_2^\vee/\cE_1).
\]
%\mnote{WZ: Why not normalize further by $q^{ds}$?}

Next we come to the geometric side. We have a natural partition
$$
(X')^r=\coprod_{\mu\in \{\pm 1\}^r} X^r.
$$
%Accordingly we write   $\cF=(\cF^{(1)},\cF^{(2)})\in \Bun_{U(n)}$ where $\cF^{(2)}=\ul{\Hom}(\cF^{(1)},\cO_X)$.\mnote{WZ: Serre dual?}
The moduli of hermitian shtukas $\Sht_{U(n)}^r$ defined in \S\ref{sec: unitary shtukas} is then partitioned into 
$$
\Sht_{U(n)}^r=\coprod_{\mu\in \{\pm 1\}^r} \Sht_{U(n)}^{\mu},
$$
and there is a natural isomorphism
$$
 \Sht_{U(n)}^{\mu}\simeq \Sht_{\GL_n}^{\mu}.
$$
Here we recall that $ \Sht_{\GL_n}^{\mu}$ is the moduli of shtukas for $\GL_{n}$, cf. \cite{YZ}, whose
$S$-points are given by the groupoid of the following data: 
\begin{enumerate}
\item $x_i \in X(S)$ for $i = 1, \ldots, r$.
\item  $\Cal{F}_0, \ldots, \Cal{F}_n\in \Bun_{\GL_n} ( S)$.
\item  An elementary modification $f_i \co \Cal{F}_{i-1} \dashrightarrow
 \Cal{F}_i$ at the graph of $x_i$, which is of upper of length 1 if $\mu_i=+1$ and of lower of length 1 if $\mu_i=-1$.%	\mnote{WZ: \cite{YZ}'s convention: upper of length 1 if $\mu_i=+1$ ?} 

\item An isomorphism $\varphi \co \Cal{F}_r \cong \ft \Cal{F}_0$. 
\end{enumerate}
In particular, $\Sht_{\GL_n}^{\mu}$ is empty unless $\sum_{i=1}^r \mu_i=0$.


For the special cycle $\cZ_\cE^r $ (cf. Definition \ref{def: Z}) associated to $\cE=(\cE_1,\cE_2)$ above, we have a partition
$$
\cZ_\cE^r =\coprod_{\mu\in \{\pm 1\}^r} \cZ^\mu_\cE,
$$
where an object in $\cZ^\mu_\cE(S)$ is an object as above in $ \Sht_{\GL_n}^{\mu}(S)$ together with maps
\begin{equation}\label{eq: t1t2}
\source{higher-siegel-weil-unitary-nonsingular-terms}
\Cal{E} _1\boxtimes \cO_S \xrightarrow{t_i^{(1)}} \Cal{F}_i  \xrightarrow{t_i^{(2)}}  \Cal{E}_2^{\vee} \boxtimes \cO_S,\quad i = 1, \ldots, r,
\end{equation}
 such that the diagram commutes
\[
\begin{tikzcd}
\Cal{E}_1 \boxtimes \cO_S \ar[d, "t_0^{(1)}"] \ar[r, equals] &  \Cal{E}_1 \boxtimes \cO_S \ar[r, equals]  \ar[d, "t_1^{(1)}"]& \ldots \ar[d] \ar[r, equals] & \Cal{E}_1\boxtimes \cO_S  \ar[r, "\sim"] \ar[d, "t_{r}^{(1)}"] & \ft (\Cal{E}_1 \boxtimes \cO_S) \ar[d, "\ft t_0^{(1)}"] \\
\Cal{F}_0  \ar[d, "t_0^{(2)}"] \ar[r, dashed, "f_0"] & \Cal{F}_1 \ar[r, dashed, "f_1"]   \ar[d, "t_1^{(2)}"]&  \ldots \ar[d] \ar[r, dashed, "f_r"] &  \Cal{F}_{r}\ar[d, "t_{r}^{(2)}"]  \ar[r, "\sim"] & \ft \Cal{F}_0  \ar[d, "\ft t_{0}^{(2)}"]\\
\Cal{E}_2^\vee \boxtimes \cO_S \ar[r, equals] &  \Cal{E}_2^\vee  \boxtimes \cO_S \ar[r, equals] & \ldots  \ar[r, equals] & \Cal{E}_2^\vee  \boxtimes \cO_S\ar[r, "\sim"] & \ft ( \Cal{E}_2^\vee \boxtimes \cO_S)
\end{tikzcd}
\]
Let $a:\cE_1\to\cE_2^\vee$ be a map of $\cO_X$-modules.  Then $\cZ^\mu_\cE(a)$ is the open-closed subscheme of $\cZ^\mu_\cE$ such that the common composition \eqref{eq: t1t2} is equal to $a \boxtimes \Id_{\cO_S} $.


For an injective $a:\cE_1\to\cE_2^\vee$, our \S\ref{s:int prob} shows that $\cZ_{\cE}^\mu(a)$ is proper over $\Spec k$ and  defines a class $[\cZ_{\cE}^\mu(a)]\in \Ch_0( \cZ_{\cE}^\mu(a))$ for each $\mu\in \{\pm 1\}^r$.  Then our main Theorem asserts 
\begin{equation}\label{eq: SWF spl}
\source{higher-siegel-weil-unitary-nonsingular-terms}
\sum_{\mu\in \{\pm 1\}^r} \deg [\cZ_{\cE}^\mu(a)] = \frac{1}{(\log q)^r}  \left(\frac{d}{ds}\right)^r\Big|_{s=0}  \left( q^{ds} \wt{E}_{a}(m(\cE),s,\Phi) \right),
\end{equation}
where $d=-(\chi(X,\cE_1)+\chi(X,\cE_2))$. We remark that $\deg [\cZ_{\cE}^\mu(a)]$ is not independent of $\mu \in \{\pm 1\}^r$, even if we restrict our attention to those $\mu$ with $\sum_{i=1}^r \mu_i = 0$. 


  



\bibliographystyle{amsalpha}
\bibliography{Bibliography}
